%==============================================================================
% Appendix D: Mass Splitting Formulae in the Spin-2 Sector
%==============================================================================
\section{Mass Splitting Formulae in the Spin-2 Sector}
\label{app:D}

This appendix collects the residual $SO(2)$ splitting, the closed-form Hessians
for all spin-2 modes, and the phase information of the spin-2 sector.

\subsection{Residual $SO(2)$ Decomposition}

Under a positive-$\varepsilon$ squash background, $SO(3)\to SO(2)$ and the
spin-2 quintet splits as follows.

\begin{table}[H]
\centering
\begin{tabular}{lll}
\toprule
Mode & Basis & Residual charge \\
\midrule
Singlet  & $T_1$          & $m=0$ \\
Doublet  & $(T_2, T_3)$   & $m=\pm 2$ \\
Doublet  & $(T_4, T_5)$   & $m=\pm 1$ \\
\bottomrule
\end{tabular}
\caption{Residual $SO(2)$ decomposition of the spin-2 quintet.}
\end{table}

\subsection{Isotropic Point}

At the isotropic point Schur's lemma gives
\begin{equation}
  H_{AB} = 2\lambda\,\delta_{AB},
  \qquad
  \mu_q^2 = 48\pi^2,
\end{equation}
with all five components fully degenerate.

\subsection{Local Splitting around the Isotropic Point}

On the axisymmetric slice $s=0$ the three representative branches are tracked as
\begin{equation}
  \mu^2_{m=0} \equiv \mu^2_\varepsilon,
  \qquad
  \mu^2_{m=\pm 2} \equiv \mu^2_s,
  \qquad
  \mu^2_{m=\pm 1} \equiv \mu^2_\delta.
\end{equation}
From a numerical scan near the isotropic point, the local linear splittings are
approximately
\begin{align}
  \mu^2_{m=0}(\varepsilon)   &\approx 48\pi^2 - 1.37\times10^3\,\varepsilon + O(\varepsilon^2), \\
  \mu^2_{m=\pm2}(\varepsilon) &\approx 48\pi^2 + 7.95\times10^2\,\varepsilon + O(\varepsilon^2), \\
  \mu^2_{m=\pm1}(\varepsilon) &\approx 48\pi^2 - 3.0\times10^2\,\varepsilon  + O(\varepsilon^2).
\end{align}
The coefficients are local slopes obtained from numerical data at $|\varepsilon|
\lesssim 0.1$ and serve as a practical summary of the branch ordering near the
isotropic point, not as exact closed-form results.

The direction of splitting is:
\begin{itemize}
  \item For positive $\varepsilon$, the $m=0$ singlet softens most rapidly.
  \item The $m=\pm 2$ doublet hardens in the same region and plays no role in the
        positive-side instability.
  \item The $m=\pm 1$ doublet lies between them and remains massive near the
        critical point.
\end{itemize}
This is consistent with the description in \S\ref{sec:spin2}.

\subsection{Closed-Form Hessians for All Spin-2 Modes}

The following closed-form expressions for $H_{ii}^q$ of the three representative
modes use the parameter $a=(1+\varepsilon)^{1/3}$ (physical parameters: $\eta=L
=\kappa=1$, $\alpha=\theta_\NY=0$, MX mode).

\paragraph{$T_1$ ($m=0$ squash singlet):}
\begin{equation}
  H_{11}^q(a,r) = \frac{-160\pi^2 a^6 r + 352\pi^2 r}{3a^8},
  \qquad a = (1+\varepsilon)^{1/3}.
\end{equation}

\paragraph{$T_2$ ($m=\pm 2$ shear doublet):}
\begin{equation}
  H_{22}^q(a,r) = \frac{128\pi^2 a^6 r - 64\pi^2 r}{a^2},
  \qquad a = (1+\varepsilon)^{1/3}.
\end{equation}

\paragraph{$T_4$ ($m=\pm 1$ off-diagonal doublet):}
\begin{equation}
  H_{44}^q(f,r) = \frac{32\pi^2 f^{12} r - 48\pi^2 f^6 r + 32\pi^2 f^3 r
                        + 48\pi^2 r}{f^8},
  \qquad f = (1+\varepsilon)^{2/3} = a^2.
\end{equation}

At the isotropic point $\varepsilon=0$ ($a=f=1$), Schur's lemma gives
\begin{equation}
  H_{11}^q = H_{22}^q = H_{44}^q = 64\pi^2 r_0,
\end{equation}
consistent with \S\ref{sec:spin2} and confirmed numerically by the scripts.

The exact critical points (both independent of $r_0$, derived from the
closed-form expressions):
\begin{align}
  H_{22}^q = 0 &\;\iff\; 2a^6 = 1
  \;\Rightarrow\;
  \varepsilon_{c-} = \frac{1}{\sqrt{2}} - 1 \approx -0.292893 \quad\text{(exact)}, \\
  H_{11}^q = 0 &\;\iff\; 5a^6 = 11
  \;\Rightarrow\;
  \varepsilon_{c+} = \sqrt{\frac{11}{5}} - 1 \approx 0.48324 \quad\text{(exact)}.
\end{align}

Script references: \texttt{paper03/spin2\_diagonal\_hessians.py}
($H_{11}$, $H_{22}$ closed forms);
\texttt{paper03/spin2\_offdiag\_hessian.py}
($H_{44}$ closed form; Method~2 rationalisation).

\subsection{Representative Spectrum Values}

\begin{table}[H]
\centering
\begin{tabular}{rrrr}
\toprule
$\varepsilon$ & $\mu^2_{m=0}$ & $\mu^2_{m=\pm 2}$ & $\mu^2_{m=\pm 1}$ \\
\midrule
$0.00$ & 473.74 & 473.74 & 473.74 \\
$0.10$ & 336.83 & 589.55 & 442.41 \\
$0.20$ & 219.78 & 685.05 & 413.32 \\
$0.40$ & 44.91  & 842.19 & 359.80 \\
$0.482$ & 0.49  & 891.77 & 341.58 \\
\bottomrule
\end{tabular}
\caption{Representative spin-2 branch values on the axiymmetric slice $s=0$.
         The soft branch is exclusively $m=0$; the $m=\pm 2$ branch moves
         separately toward the negative-side transition.}
\label{tab:spin2_spectrum}
\end{table}

\subsection{Positive Side}

On the positive side the $T_1$ singlet becomes soft first, at
\begin{equation}
  \varepsilon_{c+} = \sqrt{\tfrac{11}{5}} - 1 \approx 0.48324 \quad\text{(exact)},
  \qquad
  r_0(\varepsilon_{c+}) \approx 2.8114,
\end{equation}
where $\mu_\varepsilon^2\to 0$.  This threshold is exact and independent of $r_0$;
it is a softening threshold, not the establishment of a fully condensed phase.

\subsection{Negative Side}

On the negative side the $m=\pm 2$ sector drives the transition to the $s\neq 0$
phase.  From the closed-form expression the exact critical point is
\begin{equation}
  \varepsilon_{c-} = \tfrac{1}{\sqrt{2}} - 1 \approx -0.292893 \quad\text{(exact)},
\end{equation}
consistent with the three-indicator numerical estimate $-0.29309\pm 0.00019$.
The four-th coefficient in the Landau expansion satisfies
$b(\varepsilon)\in[2288,7805]>0$, confirming a second-order transition.

\subsection{Order Parameter Fit}

The order parameter is fitted as
\begin{equation}
  s_0 \propto |\varepsilon - \varepsilon_c|^\beta,
  \qquad
  \beta = 0.4678,\quad R^2 = 0.987,
\end{equation}
consistent with a continuous transition.  $\beta$ falls slightly below the
mean-field value $1/2$ by $1/2-\beta\approx 0.032$; this paper does not claim
to extract the universality class from this minisuperspace EFT.

\subsection{Interpretation}

The final picture of the spin-2 sector should be understood as mass splitting
organised by residual representation, not by the single squash variable.
This paper employs representation-theoretic language as the primary description.

\subsection{Computational Verification}

\begin{itemize}
  \item Isotropic point Schur degeneracy ($\mu_q^2=48\pi^2$):
        \texttt{paper03/spin2\_schur\_degeneracy.py}
  \item $SO(2)$ splitting, softening, phase transition
        (three-indicator analysis, $\varepsilon_{c\pm}$ estimation):
        \texttt{paper03/spin2\_spectrum\_splitting.py}
  \item $H_{11}^q$, $H_{22}^q$ closed forms ($\varepsilon_{c\pm}$ exact):
        \texttt{paper03/spin2\_diagonal\_hessians.py}
  \item $H_{44}^q$ closed form (Method~2 rationalisation):
        \texttt{paper03/spin2\_offdiag\_hessian.py}
\end{itemize}
